% Part: proof-theory
% Chapter: sequent-calculus
% Section: proof-examples

\documentclass[../../../include/open-logic-section]{subfiles}

\begin{document}

\olfileid{pt}{seq}{ex}

\olsection{প্রমাণের উদাহরণ}

সিকোয়েন্টের !!{proof}s নির্মাণে সাধারণত অন্তিম
সিকোয়েন্ট থেকে ওপরের দিকে এগোনো সুবিধাজনক:
সেখানে !!a{formula} সক্রিয় !!{formula} হিসেবে
বেছে নিয়ে সিকোয়েন্টটির ওপরে সংশ্লিষ্ট বিধির
পূর্বধারণা লিখি; স্বতঃসিদ্ধে পৌঁছানো পর্যন্ত
প্রক্রিয়াটি পুনরাবৃত্তি করি। যেমন, ধরি আমরা
নিচের সিকোয়েন্টটি প্রমাণ করতে চাই:
\[(!C \land !D) \lif !E \Sequent !C \lif (!D \lif !E)\]
$!C \lif (!D \lif !E)$ দেখি: এটি $!A \lif !B$
রূপের এবং সিকোয়েন্টের ডান দিকে আছে।
সম্পূর্ণ সিকোয়েন্টটি \Log{G1c}-র
$\RightR{\lif}$ বিধির সঙ্গে মেলাই:
\[\Axiom$ !A, \Gamma \fCenter \Delta, !B$
\RightLabel{\RightR{\lif}}
\UnaryInf$ \Gamma \fCenter \Delta, !A \lif !B $
\DisplayProof\]
তাহলে $\Gamma = \{(!C \land !D) \lif !E\}$,
$\Delta = \emptyset$, $!A \ident !C$ এবং
$!B \ident !D \lif !E$ নিতে হবে।
সুতরাং প্রমাণের শেষ অনুমানটি এমন হতে পারে:
\[\Axiom$ !C, (!C \land !D) \lif !E \fCenter !D \lif !E$
\RightLabel{\RightR{\lif}}
\UnaryInf$(!C \land !D) \lif !E \fCenter !C \lif (!D \lif !E)$
\DisplayProof\]
একই পদ্ধতি আবার প্রয়োগ করে
$!D \lif !E$-কে সক্রিয় সূত্র ধরলে পাই:
\[\Axiom$ !D, !C, (!C \land !D) \lif !E \fCenter !E$
\RightLabel{\RightR{\lif}}
\UnaryInf$!C, (!C \land !D) \lif !E \fCenter !D \lif !E$
\RightLabel{\RightR{\lif}}
\UnaryInf$(!C \land !D) \lif !E \fCenter !C \lif (!D \lif !E)$
\DisplayProof\]
এবার $(!C \land !D) \lif !E$-র ওপর
মনোযোগ দিয়ে \Log{G1c}-র $\LeftR{\lif}$
বিধি প্রয়োগ করতে হবে:
\[\Axiom$ \Gamma \fCenter \Delta, !A$
\Axiom$ !B, \Gamma \fCenter \Delta$
\RightLabel{\LeftR{\lif}}
\BinaryInf$ !A \lif !B, \Gamma \fCenter \Delta$
\DisplayProof\]
এতে পাই:
\[\Axiom$!D, !C \fCenter !E, !C \land !D$
\Axiom$!D, !C, !E \fCenter !E$
\RightLabel{\LeftR{\lif}}
\BinaryInf$ !D, !C, (!C \land !D) \lif !E \fCenter !E$
\RightLabel{\RightR{\lif}}
\UnaryInf$!C, (!C \land !D) \lif !E \fCenter !D \lif !E$
\RightLabel{\RightR{\lif}}
\UnaryInf$(!C \land !D) \lif !E \fCenter !C \lif (!D \lif !E)$
\DisplayProof\]
বাঁ দিকের শীর্ষ সিকোয়েন্টে
$!C \land !D$ !!{formula}টিকে
মুখ্য ধরে $\RightR{\land}$ বিধি
প্রয়োগ করলে পাই:
\[
\Axiom$!D, !C \fCenter !E, !C$
\Axiom$!D, !C \fCenter !E, !D$
\RightLabel{\RightR{\land}}
\BinaryInf$!D, !C \fCenter !E, !C \land !D$
\Axiom$!D, !C, !E \fCenter !E$
\RightLabel{\LeftR{\lif}}
\BinaryInf$ !D, !C, (!C \land !D) \lif !E \fCenter !E$
\RightLabel{\RightR{\lif}}
\UnaryInf$!C, (!C \land !D) \lif !E \fCenter !D \lif !E$
\RightLabel{\RightR{\lif}}
\UnaryInf$(!C \land !D) \lif !E \fCenter !C \lif (!D \lif !E)$
\DisplayProof
\]
এ পর্যন্ত পদ্ধতিটি \Log{G3c}-তেও চলে:
এখানে ব্যবহৃত বিধিগুলি \Log{G1c}
ও~\Log{G3c}-তে একই। এখন
!!a{proof}-এর এমন একটি খণ্ড
পেয়েছি, যার প্রতিটি শীর্ষ
সিকোয়েন্টে বাঁ ও ডান দিকে
একই !!{formula} আছে। তবে
এগুলি এখনও ঠিক স্বতঃসিদ্ধ নয়।
% BN-SRC-636 TARGET-CORRECTION-BEGIN
\par\smallskip\noindent\textnormal{\small\textbf{উৎসসংশোধন (BN-SRC-636):} প্রথম \(\RightR{\lif}\) মিলনে উৎসের \(\Gamma\) ভুল করে ডান পাশের লক্ষ্যসূত্র। সিদ্ধান্তের বাঁ পাশে থাকা একমাত্র সূত্র \((C\land D)\lif E\) দিয়ে \(\Gamma\) মেলানো হয়েছে।}
% BN-SRC-636 TARGET-CORRECTION-END
% BN-SRC-637 TARGET-CORRECTION-BEGIN
\par\smallskip\noindent\textnormal{\small\textbf{উৎসসংশোধন (BN-SRC-637):} প্রসারিত বৃক্ষের দুই পূর্বধারণাযুক্ত \(\lif\)-বাঁ ধাপটি উৎসে \(\RightR{\lif}\) লেখা; G1c/G3c সারণি ও আগের বৃক্ষ অনুযায়ী \(\LeftR{\lif}\) করা হয়েছে।}
% BN-SRC-637 TARGET-CORRECTION-END

\Log{G1c}-তে শীর্ষ সিকোয়েন্টগুলির
!!{proof}s $!A \Sequent !A$
রূপের স্বতঃসিদ্ধ থেকে দিতে হবে।
দুর্বলীকরণ-বিধি দিয়ে তা সহজেই
করা যায়। যেমন, একেবারে বাঁ দিকের
শীর্ষ সিকোয়েন্ট $!D, !C \Sequent
!E, !C$ এভাবে !!{prove} করা যায়:
\[
\Axiom$!C \fCenter !C$
\RightLabel{\LeftR{\Weakening}}
\UnaryInf$!D, !C \fCenter !C$
\RightLabel{\RightR{\Weakening}}
\UnaryInf$!D, !C \fCenter !E, !C$
\]
সব মিলিয়ে \Log{G1c}-তে
নিচের !!a{proof} পাই:
\[
\Axiom$!C \fCenter !C$
\RightLabel{\LeftR{\Weakening}}
\UnaryInf$!D, !C \fCenter !C$
\RightLabel{\RightR{\Weakening}}
\UnaryInf$!D, !C \fCenter !E, !C$
\Axiom$!D \fCenter !D$
\RightLabel{\LeftR{\Weakening}}
\UnaryInf$!D, !C \fCenter !D$
\RightLabel{\RightR{\Weakening}}
\UnaryInf$!D, !C \fCenter !E, !D$
\RightLabel{\RightR{\land}}
\BinaryInf$!D, !C \fCenter !E, !C \land !D$
\Axiom$!E \fCenter !E$
\RightLabel{\LeftR{\Weakening}}
\UnaryInf$!C, !E \fCenter !E$
\RightLabel{\LeftR{\Weakening}}
\UnaryInf$!D, !C, !E \fCenter !E$
\RightLabel{\LeftR{\lif}}
\BinaryInf$ !D, !C, (!C \land !D) \lif !E \fCenter !E$
\RightLabel{\RightR{\lif}}
\UnaryInf$!C, (!C \land !D) \lif !E \fCenter !D \lif !E$
\RightLabel{\RightR{\lif}}
\UnaryInf$(!C \land !D) \lif !E \fCenter !C \lif (!D \lif !E)$
\DisplayProof
\]
এই প্রমাণের গঠন, বিশেষত তার
উচ্চতা, $!C$, $!D$ ও~$!E$
!!{formula}s ঠিক কী তার ওপর
নির্ভর করে না। ওপরের
!!{proof}-এ এই চিহ্নগুলির
জায়গায় যেকোনো !!{formula}s
বসালেও ফলটি \Log{G1c}-র
!!{proof} থাকবে। তাই
\Log{G1c}-র প্রমাণটি
\emph{ছকভিত্তিক}।
% BN-SRC-638 TARGET-CORRECTION-BEGIN
\par\smallskip\noindent\textnormal{\small\textbf{উৎসসংশোধন (BN-SRC-638):} পূর্ণ G1c বৃক্ষের \(E\Sequent E\) থেকে বাঁ-দুর্বলীকরণে উৎসে ভুল \(C,E\Sequent C\) এসেছে; সিদ্ধান্তসূত্র অক্ষত রেখে \(C,E\Sequent E\) করা হয়েছে।}
% BN-SRC-638 TARGET-CORRECTION-END

\Log{G3c}-তে স্বতঃসিদ্ধগুলি
$!A, \Gamma \Sequent \Delta, !A$
রূপের সিকোয়েন্ট, যেখানে
$!A$ অবশ্যই আণবিক। তাই
$!D, !C \Sequent !E, !C$
\Log{G3c}-র স্বতঃসিদ্ধের
মতো দেখালেও ($\Gamma = \{!D\}$;
$\Delta = \{!E\}$), $!C$
আণবিক হলেই কেবল এটি
\emph{সত্যিই} স্বতঃসিদ্ধ।
আমাদের প্রমাণখণ্ডটি \Log{G3c}-তে
পূর্ণ প্রমাণ হবে কেবল যদি
$!C$, $!D$ ও $!E$
আণবিক !!{formula}s হয়।

কঠোর অর্থে আমরা এখনও দেখাইনি যে
\[\Log{G3c} \Proves (!C \land !D) \lif !E \fCenter !C \lif (!D \lif
!E),\]
তবু বাস্তবে এটুকুই করা প্রয়োজন
(এবং এটুকুই করা যায়)। কারণ
$!A$ নিজে আণবিক না হলেও
$!A, \Gamma \Sequent \Delta, !A$
রূপের প্রতিটি সিকোয়েন্ট
\Log{G3c}-তে প্রমাণযোগ্য।
$!A$-র গভীরতা~$\depth{!A}$-র
ওপর আরোহ করে এটি দেখাতে পারি।
% BN-SRC-639 TARGET-CORRECTION-BEGIN
\par\smallskip\noindent\textnormal{\small\textbf{উৎসসংশোধন (BN-SRC-639):} আগের অনুচ্ছেদে G1c-র পূর্ণ প্রমাণ ইতিমধ্যে দেখানো হয়েছে; আণবিক নয় এমন শীর্ষসূত্রের জন্য অসম্পূর্ণ খণ্ডটি G3c-র। অতএব পরের প্রদর্শিত প্রমাণযোগ্যতার দাবিতে উৎসের G1c-র বদলে G3c।}
% BN-SRC-639 TARGET-CORRECTION-END

\begin{defn}\ollabel{defn:depth}
!!a{formula}-র \emph{গভীরতা}~$\depth{!A}$
নিচের মতো সংজ্ঞায়িত:
\begin{enumerate}
  \item $!A$ আণবিক হলে $\depth{!A} = 0$।
  \item $!A \ident \lnot !B$, $!A \ident \lforall[x][!B(x)]$,
  অথবা $!A \ident \lexists[x][!B(x)]$ হলে
  $\depth{!A} = \depth{!B} + 1$।
  \item $!A \ident (!B \land !C)$, $(!B \lor !C)$, অথবা $(!B \lif
  !C)$ হলে $\depth{!A} = \max(\depth{!B},\depth{!C}) + 1$।
\end{enumerate}
\end{defn}
% BN-SRC-640 TARGET-CORRECTION-BEGIN
\par\smallskip\noindent\textnormal{\small\textbf{উৎসসংশোধন (BN-SRC-640):} গভীরতার সংজ্ঞা, G1c-র পরিমাণসূচক-আরোহ ও দুই অনুশীলনে উৎসের অনির্দিষ্ট \(A\), \(A(x)\), \(B(x)\) ও \(B(c)\) লেখা নিয়মসারণি ও প্রদর্শিত বৃক্ষের সঙ্গে মিলিয়ে \(!B(x)\), \(!B(c)\) করা হয়েছে। পরিমাণসূচকের ভেতরের সূত্রের গভীরতাই এক কম।}
% BN-SRC-640 TARGET-CORRECTION-END

যেকোনো পদ~$t$-এর জন্য
$\depth{A(x)} = \depth{A(t)}$
প্রমাণ করা কঠিন নয়। আমাদের
জন্য গুরুত্বপূর্ণ হলো, প্রতিটি
বিধিতে মুখ্য !!{formula}-র
গভীরতা সক্রিয় !!{formula}s-এর
গভীরতার চেয়ে বেশি। যেমন:
\begin{align*}
\depth{P(x)} = \depth{P(a)} & = 0,\\
\depth{\lforall[x][P(x)]} & = 1,\\
\depth{(\lforall[x][P(x)] \lif P(a))} & = \max(1,0) + 1 = 2.
\end{align*}
এই ধর্ম ব্যবহার করে
$\depth{!A}$-র ওপর
আরোহে নানা ফল প্রমাণ
করতে পারি; যেমন:

\begin{prop}\ollabel{prop:G1c-axioms}
যেকোনো~$!A$-র জন্য
$!A \Sequent !A$-র
$\Log{G1c}$-তে এমন !!a{proof}
আছে, যার সব স্বতঃসিদ্ধ আণবিক।
\end{prop}

\begin{proof}
  $\depth{!A}$-র ওপর আরোহ করি।
  $\depth{!A} = 0$ হলে $!A$
  আণবিক, আর $!A \Sequent !A$
  নিজেই চাওয়া রূপের
  \Log{G1c}-!!{proof}।
  $\depth{!A} > 0$ হলে
  $!A$ কম গভীরতার
  !!{formula}s থেকে গঠিত;
  যেমন $!A \ident \lnot !B$,
  $!A \ident (!B \land !C)$,
  অথবা $!A \ident \lforall[x][!B(x)]$।
  আরোহের অনুমান: $\depth{!D} <
  \depth{!A}$ এমন প্রতিটি $!D$-র
  জন্য আণবিক স্বতঃসিদ্ধ থেকে
  $\Log{G1c} \Proves !D \Sequent !D$।

  $!A \ident \lnot !B$ হলে
  $\depth{!B} < \depth{!A}$।
  আরোহের অনুমানে আণবিক
  স্বতঃসিদ্ধ থেকে $!B \Sequent !B$-র
  \Log{G1c}-!!a{proof}~$\pi_1$
  আছে। এটিকে নিচের মতো
  $\lnot !B \Sequent \lnot !B$-র
  !!a{proof} পর্যন্ত বাড়াই:
  \[
  \AxiomC{}
  \RightLabel{$\pi_1$}
  \Deduce$!B \fCenter !B$
  \RightLabel{\LeftR{\lnot}}
  \UnaryInf$\lnot !B, !B \fCenter $
  \RightLabel{\RightR{\lnot}}
  \UnaryInf$\lnot !B \fCenter \lnot !B$
  \DisplayProof
\]

  $!A \ident (!B \land !C)$ হলে
  $\depth{!B} < \depth{!A}$
  এবং $\depth{!C} < \depth{!A}$।
  আরোহের অনুমানে আণবিক
  স্বতঃসিদ্ধ থেকে $!B
  \Sequent !B$ ও
  $!C \Sequent !C$-র
  \Log{G1c}-!!{proof}s~$\pi_1$
  ও~$\pi_2$ আছে।
  এগুলিকে নিচের মতো
  $!B \land !C \Sequent
  !B \land !C$-র
  !!a{proof} পর্যন্ত বাড়াই:
  \[
  \AxiomC{}
  \RightLabel{$\pi_1$}
  \Deduce$!B \fCenter !B$
  \RightLabel{\LeftR{\land}}
  \UnaryInf$!B \land !C \fCenter !B$
  \AxiomC{}
  \RightLabel{$\pi_2$}
  \Deduce$!C \fCenter !C$
  \RightLabel{\LeftR{\land}}
  \UnaryInf$!B \land !C \fCenter !C$
  \RightLabel{\RightR{\land}}
  \BinaryInf$!B \land !C \fCenter !B \land !C$
  \DisplayProof
  \]
% BN-SRC-641 TARGET-CORRECTION-BEGIN
\par\smallskip\noindent\textnormal{\small\textbf{উৎসসংশোধন (BN-SRC-641):} G1c-র \(\land\)-পরিচয় প্রমাণের উৎসবৃক্ষে প্রথম \(\pi_1\) সিদ্ধান্তে বাড়তি কমা, বাঁ-প্রসঙ্গে অযৌক্তিক পুনরাবৃত্তি এবং শেষে ভুল \(B\land C\Sequent B\) ছিল। G1c-র দুটি বিকল্প \(\LeftR{\land}\) ও একটি \(\RightR{\land}\) দিয়ে একই দুই আরোহী প্রমাণ থেকে বৈধ \(B\land C\Sequent B\land C\) বৃক্ষ দেওয়া হয়েছে।}
% BN-SRC-641 TARGET-CORRECTION-END

  এবার ধরি $!A \ident \lforall[x][!B(x)]$।
  $\lforall[x][!B(x)]$-এ নেই এমন
  !!a{constant}~$c$ নিই।
  তাহলে $\depth{!B(c)} =
  \depth{!B(x)} < \depth{!A}$।
  আরোহের অনুমানে আণবিক
  স্বতঃসিদ্ধ থেকে $!B(c)
  \Sequent !B(c)$-র
  \Log{G1c}-!!a{proof}~$\pi_1$
  আছে। এটিকে নিচের মতো
  $\lforall[x][!B(x)] \Sequent
  \lforall[x][!B(x)]$-র
  !!a{proof} পর্যন্ত বাড়াই:
  \[
  \AxiomC{}
  \RightLabel{$\pi_1$}
  \Deduce$!B(c) \fCenter !B(c)$
  \RightLabel{\LeftR{\lforall}}
  \UnaryInf$\lforall[x][!B(x)] \fCenter !B(c)$
  \RightLabel{\RightR{\lforall}}
  \UnaryInf$\lforall[x][!B(x)] \fCenter \lforall[x][!B(x)]$
  \DisplayProof
\]

  বাকি ক্ষেত্রগুলি অনুশীলন
  হিসেবে থাক।
\end{proof}

\begin{prob}
  $!A \ident (!B \lor !C)$,
  $!A \ident (!B \lif !C)$
  এবং $!A \ident \lexists[x][!B(x)]$
  ক্ষেত্রগুলি দেখিয়ে
  \olref{prop:G1c-axioms}-এর
  প্রমাণ সম্পূর্ণ করো।
\end{prob}

আমাদের উদ্দেশ্যে নিচের
!!{proof}টিও ব্যবহার
করতে পারতাম:
\[
\AxiomC{}
\RightLabel{$\pi_1$}
\Deduce$!B \fCenter !B$
\RightLabel{\RightR{\lnot}}
\UnaryInf$\fCenter !B, \lnot !B$
\RightLabel{\LeftR{\lnot}}
\UnaryInf$\lnot !B \fCenter \lnot !B$
\DisplayProof
\]
$!A \ident \lnot !B$
ক্ষেত্রে। কিন্তু এমন সিকোয়েন্ট
ব্যবস্থা আছে যেখানে এটি চলত না।
স্বজ্ঞাবাদী সিকোয়েন্ট কলন
\Log{G1i}, \Log{G1c}-র মতোই;
পার্থক্য হলো, প্রতিটি সিকোয়েন্টের
ডান দিকে সর্বাধিক একটি
!!{formula} থাকতে পারে।
$\Delta = \emptyset$ হলে
প্রথম !!{proof}টি
\Log{G1i}-তেও !!a{proof}
হতো, কিন্তু দ্বিতীয়টি নয়:
তার একটি সিকোয়েন্টের
ডান দিকে $!B$ ও
$\lnot !B$ দুটিই আছে।

লক্ষ করো, $!A \ident
\lforall[x][!B(x)]$ ক্ষেত্রেও
\Log{G1c}-তে বিধিগুলির
অন্য ক্রম নিতে পারতাম না।
নিচেরটি \emph{!!a{proof} নয়}:
  \[
  \AxiomC{}
  \RightLabel{$\pi_1$}
  \Deduce$!B(c) \fCenter !B(c)$
  \RightLabel{\RightR{\lforall}}
  \UnaryInf$!B(c) \fCenter \lforall[x][!B(x)]$
  \RightLabel{\LeftR{\lforall}}
  \UnaryInf$\lforall[x][!B(x)] \fCenter \lforall[x][!B(x)]$
  \DisplayProof
\]
কারণ এতে $\RightR{\lforall}$-এর
আইগেনচল-শর্ত লঙ্ঘিত হয়েছে।

\begin{prop}\ollabel{prop:G3c-axioms}
$!A, \Gamma \Sequent \Delta, !A$
রূপের যেকোনো সিকোয়েন্ট
\Log{G3c}-তে প্রমাণযোগ্য;
এখানে $!A$ \emph{আবশ্যিকভাবে}
আণবিক হতে হবে না।
\end{prop}

\begin{proof}
\olref{prop:G1c-axioms}-এর
প্রমাণের সঙ্গে এই প্রমাণের
অনেক মিল, তবে আরোহের
অনুমানটি সাবধানে নিতে হবে।
$n = \depth{!A} = 0$
ক্ষেত্রটি আবার সহজ।
$n=0$ হলে $!A$ আণবিক,
আর $!A, \Gamma \Sequent
\Delta, !A$ \Log{G3c}-র
স্বতঃসিদ্ধ।

এবার ধরি $n>0$।
আবার ক্ষেত্রভাগ করি।
আরোহের অনুমান: $\depth{!D}
< n$ এমন সব $!D$ এবং
যেকোনো $\Pi$ ও
$\Lambda$-র জন্য
$\Log{G3c} \Proves
!D, \Pi \Sequent
\Lambda, !D$।
$!A \ident (!B \land !C)$
ক্ষেত্রটি দেখি।
আরোহের অনুমান দুবার
প্রয়োগ করি। প্রথমে
$!D = !B$,
$\Pi = !C, \Gamma$ ও
$\Lambda = \Delta$
নিয়ে $!B, !C, \Gamma
\Sequent \Delta, !B$-র
!!a{proof}~$\pi_1$ পাই।
তারপর $!D = !C$,
$\Pi = !B, \Gamma$ ও
$\Lambda = \Delta$
নিয়ে $!B, !C, \Gamma
\Sequent \Delta, !C$-র
!!a{proof}~$\pi_2$ পাই।
$!B \land !C, \Gamma
\Sequent \Delta,
!B \land !C$-র
!!{proof}টি হলো:
\[
\AxiomC{}
\RightLabel{$\pi_1$}
\Deduce$!B, !C, \Gamma \fCenter \Delta, !B$
\RightLabel{\LeftR{\land}}
\UnaryInf$!B \land !C, \Gamma \fCenter \Delta, !B$
\AxiomC{}
\RightLabel{$\pi_2$}
\Deduce$!B, !C, \Gamma \fCenter \Delta, !C$
\RightLabel{\LeftR{\land}}
\UnaryInf$!B \land !C, \Gamma \fCenter \Delta, !C$
\RightLabel{\RightR{\land}}
\BinaryInf$!B \land !C, \Gamma \fCenter \Delta, !B \land !C$
\DisplayProof
\]
% BN-SRC-642 TARGET-CORRECTION-BEGIN
\par\smallskip\noindent\textnormal{\small\textbf{উৎসসংশোধন (BN-SRC-642):} G3c-র আরোহের অনুমানে উৎসের উল্টো $D!$ সূত্রলিখন $!D$ করা হয়েছে। $\pi_2$-র জন্য উৎসের $\Pi=B,\Lambda$ ভুল প্রসঙ্গ; বৃক্ষের $\Pi=B,\Gamma$ বসানো হয়েছে।}
% BN-SRC-642 TARGET-CORRECTION-END

$!A \ident \lforall[x][!B(x)]$
হলে $\lforall[x][!B(x)]$,
$\Gamma$ কিংবা~$\Delta$-তে
নেই এমন !!a{constant}
বেছে নিই। আরোহের অনুমানে
$!D = !B(c)$,
$\Pi = \lforall[x][!B(x)], \Gamma$
ও $\Lambda = \Delta$
নিই। এতে $!B(c),
\lforall[x][!B(x)], \Gamma
\Sequent \Delta, !B(c)$-র
!!a{proof}~$\pi_1$
পাই; সেখান থেকে:
\[
\AxiomC{}
\RightLabel{$\pi_1$}
\Deduce$!B(c), \lforall[x][!B(x)], \Gamma \fCenter \Delta, !B(c)$
\RightLabel{\LeftR{\lforall}}
\UnaryInf$\lforall[x][!B(x)], \Gamma \fCenter \Delta, !B(c)$
\RightLabel{\RightR{\lforall}}
\UnaryInf$\lforall[x][!B(x)], \Gamma \fCenter \Delta, \lforall[x][!B(x)]$
\DisplayProof
\]
$c$ যেভাবে বেছে নিয়েছি,
তাতে $\RightR{\lforall}$-এর
আইগেনচল-শর্ত পূরণ হয়।
\end{proof}

\begin{prob}
  $!A \ident (!B \lif !C)$
  এবং $!A \ident
  \lexists[x][!B(x)]$
  ক্ষেত্রগুলি দেখিয়ে
  \olref{prop:G3c-axioms}-এর
  প্রমাণ এগিয়ে নাও।
\end{prob}

\end{document}
